\chapter{Background and Related Work}

\section{Sign Language and Fingerspelling}

A signed language conveys meaning through handshape, orientation, location,
movement, and non-manual markers such as facial expression. ASL has a
dedicated set of one-handed handshapes for the letters of the alphabet, used
for \emph{fingerspelling}—spelling out words letter by letter—which is common
for names and for words that lack a dedicated sign. Most letters are static
handshapes; two, \textbf{J} and \textbf{Z}, are defined by a traced motion.
Word-level signs are typically dynamic and may be one- or two-handed. Because
the device in this project instruments a single hand, the vocabulary is
restricted to fingerspelling and to one-handed word signs.

A practical consequence for sensing is that several fingerspelling letters
differ only in subtle thumb placement against a closed fist—for example
\textbf{A}, \textbf{S}, \textbf{T}, \textbf{M}, and \textbf{N}. The ability of
any sensing modality to separate these handshapes is therefore a key indicator
of its resolution.

\section{Approaches to Sign Language Recognition}

Sign-language recognition (SLR) systems fall broadly into two families.

\textbf{Vision-based} systems use one or more cameras and computer-vision or
deep-learning models to interpret hand and body pose. They require no
instrumentation of the signer but depend on controlled lighting, a clear line
of sight, and significant computation, and they raise privacy concerns because
they continuously capture video.

\textbf{Sensor-based (glove)} systems instrument the hand directly with
sensors such as flex sensors, which measure finger bend, or inertial
measurement units (IMUs), which measure orientation and motion. Gloves work in
any lighting, do not capture the surroundings, and can run on small
microcontrollers, at the cost of requiring the user to wear the device. This
project adopts the glove approach using IMUs, which—unlike flex sensors—capture
both static finger orientation and dynamic motion, and are available as
inexpensive, easily interfaced modules.

\section{Inertial Measurement Units and Sensor Fusion}

An IMU such as the InvenSense MPU6050 \cite{mpu6050} integrates a three-axis
gyroscope and a three-axis accelerometer. The accelerometer measures the
gravity vector (giving an absolute tilt reference) plus linear acceleration;
the gyroscope measures angular rate. Neither sensor alone yields a stable
orientation: accelerometer-only tilt is corrupted by motion, while integrating
the gyroscope alone accumulates drift.

\emph{Sensor fusion} combines the two to obtain a stable estimate. The Kalman
filter \cite{kalman} is the classical optimal estimator for this purpose and is
widely used to fuse accelerometer and gyroscope measurements into pitch and
roll angles. A practical complication is that the gyroscope's zero-rate bias
drifts with temperature; this is commonly mitigated by continuously
re-estimating the bias whenever the device is detected to be still. Both
techniques are used in this project's firmware (Chapter~4).

A fundamental limitation of a six-axis IMU is that yaw (rotation about the
gravity axis) cannot be referenced absolutely without a magnetometer, and that
small static orientation differences are near the noise floor of a low-cost
part. These limitations directly shape the feature-representation findings of
Chapter~5.

\section{Wireless Protocols for Wearables}

Body-worn sensor networks require a short-range, low-power link. Two relevant
technologies are used here. \textbf{ESP-NOW} \cite{espnow} is a connectionless,
low-latency protocol built on the Wi-Fi physical layer by Espressif, well
suited to many-to-one streaming between ESP32-family devices without the
overhead of a full Wi-Fi association. \textbf{Bluetooth Low Energy (BLE)} is
the standard for communicating with smartphones; its GATT profile exposes
\emph{characteristics} that a phone can subscribe to for streaming
notifications. The system in this thesis uses ESP-NOW for finger-to-finger
communication and BLE for the final hop to the phone (Chapter~3).

\section{Machine Learning at the Edge}

Running inference on a constrained device (``edge'' or ``TinyML'') requires
small models and efficient execution. Two ideas are central. First,
\textbf{one-dimensional convolutional neural networks (1-D CNNs)} are an
effective and lightweight architecture for time-series and human-activity
recognition \cite{ronao2016, hatcnn}, learning local temporal patterns with far
fewer parameters than recurrent models. Second, \textbf{quantisation} converts
a trained 32-bit floating-point model to 8-bit integers, shrinking it roughly
four-fold and accelerating inference, usually with negligible accuracy loss;
TensorFlow Lite \cite{tflite} provides this conversion and a portable runtime.
The models in this project are 1-D CNNs trained with the Edge Impulse platform
\cite{edgeimpulse} and deployed as int8 TensorFlow Lite files inside a mobile
application.
